\section{Focused covariants of the four-qubit sextic}\label{sec:sextic-general}

Throughout this section $n$ is the number of qubits and $N=\|z\|^2$ (in the four-qubit
sections the norm squared is denoted by $N$, and in the seven-qubit entropy certificate by $\nu$), and $W_1=\sum_i\|\mathbf r_i\|^2$ with
$\mathbf r_i^a=\langle z|\sigma_a^{(i)}|z\rangle$.
The four-qubit moment identity \eqref{eq:fqe-remainder} is not an ordinary positive
three-copy replica operator: the latter has a negative sector on the
copy-symmetric cube (Appendix~\ref{sec:fqe-correlation}).
The successful proof uses the spin-flipped mixed tensor
$w=z_x\otimes z_y\otimes\widetilde z_t$,
$\widetilde z=\epsilon^{\otimes n}\bar z$, with
$\epsilon=\left(\begin{smallmatrix}0&1\\-1&0\end{smallmatrix}\right)$.
A complementary Pauli-coordinate description of third moments is derived in
\eqref{eq:pauli-triple} below.
At one site write
\[
 L|xyt\rangle=\epsilon_{yt}|x\rangle,\qquad
 M|xyt\rangle=\epsilon_{xt}|y\rangle,
\]
and let $S_i$ project the local three-copy tensor at site $i$ onto
$\operatorname{Sym}^3(\mathbb C^2)$.

\begin{proposition}[focused-site sextic covariants]\label{prop:focus}
For a subset $I\subseteq[n]$ and $0\le r\le n-|I|$, define
\[
 \Phi_{I,r}(z)=\left(\prod_{i\in I}S_i\right)
 \sum_{\substack{A\subseteq I^c\\|A|=r}}
 \left(\prod_{j\in A}M_j\right)
 \left(\prod_{j\in I^c\setminus A}L_j\right)
 (z_x\otimes z_y\otimes\widetilde z_t).
\]
Each $\Phi_{I,r}$ is a local-$SU(2)$ covariant of bidegree $(2,1)$ (quadratic in $z$, linear in $\bar z$), and
\begin{equation}\label{eq:focus-reflection}
 \Phi_{I,r}=\Phi_{I,n-|I|-r}.
\end{equation}
Its Gram matrix $G_I(z)=(\langle\Phi_{I,r},\Phi_{I,s}\rangle)_{r,s}$
is positive semidefinite. After~\eqref{eq:focus-reflection}, at most
$\lfloor(n-|I|)/2\rfloor+1$ distinct vectors remain.

For $n=4$ and $I=\{i\}$ put $\mathcal R_i=4\Phi_{i,1}-15\Phi_{i,0}$, a residual different from the
$V_i-8U_i$ of \eqref{eq:fqe-remainder}. The
homogeneous sextic identity then has the reduced form
\begin{equation}\label{eq:focus-four}
32(17NP-12T-12N^3)
=3\sum_{i=1}^4\left(
 15\|\Phi_{i,0}\|^2+
 \|4\Phi_{i,1}-15\Phi_{i,0}\|^2\right),
\end{equation}
where $N=\|z\|^2$, $P=\sum_{j=2}^4\tr R_{1j}^2$ and
$T=\sum_{j=2}^4\tr R_{1j}^3$.
\end{proposition}
\begin{proof}
Interchanging the two identical factors $z_x,z_y$ swaps $L$ with
$M$ at every contracted site, while the symmetric projectors at the
focused sites are unchanged. This gives~\eqref{eq:focus-reflection}.
Gram positivity follows from the Hilbert-space inner product.
In the notation of Appendix~\ref{sec:fqe-sextic}, $U_i=\Phi_{i,0}$ and
$V_i=\sum_{r=0}^3\Phi_{i,r}=2(\Phi_{i,0}+\Phi_{i,1})$.
Thus $\mathcal R_i=2V_i-19U_i=2Y_i-7U_i$ with $Y_i=V_i-6U_i$. Writing $F=17NP-12T-12N^3$ for the
left side of~\eqref{eq:focus-four} divided by $32$, the contractions
\eqref{eq:fqe-cubic-mixed-norms} give
\[
 \sum_i\|\mathcal R_i\|^2=4\sum_i\|Y_i\|^2-28\operatorname{Re}\sum_i\langle U_i,Y_i\rangle+49\sum_i\|U_i\|^2
 =\tfrac{32}3F-5NW_1 .
\]
Hence $32F=15NW_1+3\sum_i\|\mathcal R_i\|^2$, and $\sum_i\|U_i\|^2=NW_1/3$ turns this
into~\eqref{eq:focus-four}. SF-1 checks the identity on one integer four-qubit tensor by exact rational arithmetic.
\end{proof}

\begin{lemma}[the exact-support Pauli selection rule]\label{lem:focus-pauli}
Let $R_I$ be the unnormalised reduction of an arbitrary pure $n$-qubit
state to the focused sites $I$. Then, up to the fixed sign
$(-1)^{n-|I|}$,
\[
 \Phi_{I,0}= \left(\prod_{i\in I}S_i\right)
 \left(z_x\otimes (R_I\epsilon_I^T)_{y_I,t_I}\right),
 \qquad \epsilon_I=\epsilon^{\otimes|I|},
\]
with the $x$-slots on $I^c$ left intact. In particular, if
$\langle z|P|z\rangle=0$ for every Pauli word $P$ of exact support
$I$, then $\Phi_{I,0}(z)=0$.
\end{lemma}
\begin{proof}
At a contracted site,
$\sum_t\epsilon_{yt}\epsilon_{tu}=-\delta_{yu}$,
so the two remaining replicas form $R_I$. Expand
$R_I=2^{-|I|}\sum_{P\in\{I,X,Y,Z\}^{I}}
\langle z|P|z\rangle P$.
If $P$ has an identity factor at some $i\in I$, then
$\operatorname{vec}(I\epsilon^T)$ is a local singlet in the $y,t$
slots. Its tensor product with any $x$-qubit is killed by $S_i$.
Only terms supported on every site survive.
\end{proof}

\begin{proposition}[complete scalar cubic gradients; analytic, SF-5 controls]
\label{prop:zero-focus}
For any $n$, let $Q_r(z)=\sum_{|S|=r}\tr R_S^2$, $N=\|z\|^2$,
and $V_r(z)=\sum_{|S|=r}(R_S\otimes I_{S^c})z$.
Then
\[
 \Phi_{\varnothing,r}(z)=(-1)^nV_r(z),\qquad
 \frac{\partial Q_r}{\partial\bar z}=2V_r(z),\qquad
 \langle z,V_r(z)\rangle=Q_r(z).
\]
More generally, let $v$ be a local-unitary scalar covariant of bidegree
$(2,1)$ with real coefficients in the $L,M$ tensor basis. Its closure
$p_v=\langle z,v(z)\rangle$ is real and
$v=\tfrac12\partial_{\bar z}p_v$.
The vectors $V_S=(R_S\otimes I)z$, with one $S$ from each complementary
pair $\{S,S^c\}$, form a basis of this real space, of dimension $2^{n-1}$.
Its permutation-invariant subspace has basis
$V_0,\ldots,V_{\lfloor n/2\rfloor}$ and dimension
$\lfloor n/2\rfloor+1$. Thus the zero-focus weight sums exhaust
the symmetric scalar cubic source.
Consequently, for any fixed real $\lambda$,
\begin{equation}\label{eq:gradient-sextic}
 \mathcal G_{r,\lambda}(z):=
 \|\Phi_{\varnothing,r}(z)-\lambda\Phi_{\varnothing,0}(z)\|^2
 =\|V_r(z)-\lambda Nz\|^2\ge0
\end{equation}
is a universal homogeneous sextic. It vanishes exactly when $z$
is a projective stationary point of $Q_r/N^2$ with
$Q_r/N^2=\lambda$, for nonzero $z$. Moreover,
\begin{equation}\label{eq:gradient-octic}
 \begin{split}
 N\mathcal G_{r,\lambda}(z)
 &=\bigl(N\|V_r\|^2-Q_r^2\bigr)+(Q_r-\lambda N^2)^2,\\
 N\|V_r\|^2-Q_r^2
 &=\sum_{a<b}|z_a(V_r)_b-z_b(V_r)_a|^2\ge0.
 \end{split}
\end{equation}
Along the seven-qubit icosahedral family $\psi_x$ ($0\le x\le1/2$),
\begin{equation}\label{eq:zero-focus-family}
\begin{aligned}
 &\Phi_{\varnothing,1}(\psi_x)=\tfrac72\Phi_{\varnothing,0}(\psi_x),\\
 &4(104x-7)\Phi_{\varnothing,3}(\psi_x)-2(524x-37)\Phi_{\varnothing,2}(\psi_x)
   +(96x^2+3628x-263)\Phi_{\varnothing,0}(\psi_x)=0.
\end{aligned}
\end{equation}
At $\Psiico=\psi_{37/524}$ the second relation simplifies to
$\Phi_{\varnothing,3}=\frac{605}{131}\Phi_{\varnothing,0}$.
(Both relations are proved below; SF-5 checks them exactly.)
In fact the entire zero-focus Gram on this family is
\begin{equation}\label{eq:zero-focus-gradient-gram}
 G(x)=q(x)q(x)^T+\frac{x(1-2x)}8q'(x)q'(x)^T,
 \qquad q=(1,\tfrac72,21\pi_2,35\pi_3)^T.
\end{equation}
It has rank two for $0<x<1/2$ and rank one at both endpoints.
For the two source combinations $V_0$ and $131V_2-52V_3$, its
$2\times2$ Gram determinant is exactly $72x(1-2x)$.
Thus $\mathcal G_{3,605/131}$ is a nontrivial
sextic SOS vanishing at $\Psiico$, while at $x=x_*$ the two relations
in~\eqref{eq:zero-focus-family} give the $\operatorname{Sym}^7\mathbb R^2$
three-copy kernel features of Theorem~\ref{thm:seven-entropy-localized-identity}.
Neither gives a lower bound on $\pi_3$; Theorem~\ref{thm:sextic-sos-obstruction}
excludes a sharp fixed-coefficient cubic SOS certificate even after including this form.
\end{proposition}
\begin{proof}
At an $L$ site, the first replica index survives and the other
two are identified after spin flip; at an $M$ site the second
replica index survives and the first is identified with the
conjugate. For a set $S$ of $M$ sites, the surviving tensor is
$(-1)^n(R_S\otimes I_{S^c})z$.
Differentiating $\tr R_S^2$ by the product rule yields
$2(R_S\otimes I)z$. Euler's homogeneous identity gives the
inner product.

To prove that these scalar cubics exhaust the possibilities, contract the output of a general scalar cubic with
$\bar z$. Locally its two-copy invariant operator is a real combination
of identity and swap. Their tensor products are Hermitian and commute
with the simultaneous copy exchange. Since the two holomorphic inputs
are already symmetric, the two antiholomorphic inputs are symmetric
as well; hence $p_v$ is real and $\partial_{\bar z}p_v=2v$.
All quartic local-unitary invariants of bidegree $(2,2)$ are linear combinations of the
swap expectations $\tr R_S^2$. On the symmetric square, complementary
swap words agree. These are their only relations: the local
antisymmetric-site sets $A$ have even cardinality, and the words act
on them by the characters $(-1)^{|S\cap A|}$. Characters of this
even-subset group are indexed exactly by $S$ modulo complementation
and are linearly independent. Each such sector occurs on qubits.
Expectations on $z^{\otimes2}$ determine a Hermitian operator on
the symmetric square, by its polynomial coefficients.
This proves the full basis statement. After site averaging, only
$|S|$ remains. The weight-$r$ swap sum has value
$[t^r](1+t)^{n-|A|}(1-t)^{|A|}$, a polynomial of degree $r$ in $|A|$
with leading coefficient $(-2)^r/r!$. For
$0\le r\le\lfloor n/2\rfloor$ these polynomials are independent on
the equally many distinct even values of $|A|$, proving the symmetric
basis statement.

Projective stationarity of $Q_r/N^2$ is
$V_r=(Q_r/N)z$, proving the equality criterion.
The octic equation is the orthogonal decomposition of
$V_r$ into its component along $z$ and its perpendicular
component; its final equality is the usual complex
Lagrange identity. Along the icosahedral family $\psi_x$,
one-body maximal mixing gives the first relation in~\eqref{eq:zero-focus-family}.
For $r\in\{2,3\}$, the fixed-space argument of
Lemma~\ref{lem:seven-family-symmetric-criticality} places the real gradient
$4V_r(\psi_x)$ in the plane
$\operatorname{span}\{U,V\}=\operatorname{span}\{\psi_x,\psi_x'\}$, where
$\psi_x'\perp\psi_x$, $\langle\psi_x,V_r(\psi_x)\rangle=Q_r(\psi_x)$ and, by the gradient formula,
$\langle\psi_x',V_r(\psi_x)\rangle=\tfrac14 Q_r'(\psi_x)$.
Hence $V_r=Q_r\psi_x+\tfrac14Q_r'\psi_x'/\|\psi_x'\|^2$.
Since $30(a')^2+2(b')^2=1/[2x(1-2x)]$, taking inner products gives
\eqref{eq:zero-focus-gradient-gram}. The explicit purity polynomials give
$131q_2'-52q_3'=24$. The displayed Gram determinant is therefore
$24^2x(1-2x)/8=72x(1-2x)$, proving rank two throughout the open
interval; continuity gives rank one at its endpoints.
Eliminating $\psi_x'$ between $r=2$ and $r=3$ gives
$Q_2'\Phi_{\varnothing,3}-Q_3'\Phi_{\varnothing,2}-(Q_2'Q_3-Q_3'Q_2)\Phi_{\varnothing,0}=0$;
substituting $Q_2=21\pi_2$ and $Q_3=35\pi_3$ from Theorem~\ref{thm:ico} yields the second
relation in~\eqref{eq:zero-focus-family}.
Thus the two displayed relations span its kernel for every interior $x$:
their source rows are independent, since $104x-7$ and $524x-37$
cannot vanish together. SF-5 compares the closed Gram formula with every
coefficient of an independent exact tensor contraction, and checks its
specialisation at $x=37/524$ with null vectors
$(-7,2,0,0)$ and $(-605,0,0,131)$ in the
$\Phi_{\varnothing,0},\ldots,\Phi_{\varnothing,3}$ basis.
\end{proof}

\begin{remark}[scope of cubic-contraction searches]
Theorem~\ref{thm:sextic-sos-obstruction} rules out every fixed-coefficient
cubic SOS certificate of the sharp seven-qubit purity bound,
including all nonempty-focus forms, zero-focus gradients and their
cross-sector couplings. The
covariant identities above remain useful algebraic descriptions;
they do not evade the complete three-copy obstruction.
\end{remark}

\begin{remark}[joint quadratic/cubic Pauli enumerators]
For a trace-one density matrix $\rho$ define $r_P=\tr(\rho P)$ for Hermitian
Pauli strings $P$. The ordinary weight enumerator
$A_j=\sum_{\operatorname{wt}P=j}r_P^2$ controls mean subsystem
purity. There is a natural signed \emph{triple} enumerator:
\[
 C_j=\!\!\sum_{\substack{P,Q:[P,Q]=0\\
           |\operatorname{supp}P\cup\operatorname{supp}Q|=j}}
 \eta(P,Q)\,r_Pr_Qr_{R(P,Q)},
 \qquad PQ=\eta(P,Q)R(P,Q),\ \eta(P,Q)\in\{\pm1\},
\]
The sum is over ordered pairs of Pauli words, including the identity;
$R(P,Q)$ is the canonical Hermitian Pauli word.
For $b_k=\binom nk^{-1}\sum_{|S|=k}\tr\rho_S^3$,
\begin{equation}\label{eq:pauli-triple}
 b_k=\frac1{4^k\binom nk}
 \sum_{j=0}^k\binom{n-j}{k-j}C_j.
\end{equation}
Expand $\rho_S^3$ in the Pauli basis: products of
anticommuting $P,Q$ cancel between orders and every commuting
product is $\pm R(P,Q)$. Group by the union of supports to obtain
the binomial coefficient. Thus $A_j,C_j$ are natural joint
second/third-moment coordinates. For pure states the third moments are sextic in the amplitudes;
unlike $A_j$, the individual $C_j$ need not be positive. This coordinate identity is independent
of the covariant certificates and is not used in the cubic-SOS obstruction below.
\end{remark}


\section{The complete cubic-SOS obstruction}
\label{sec:sextic-all-sos}

For seven qubits let $N(z)=\|z\|^2$, let $\mathcal F_3(z)$ be the homogeneous
quartic mean three-body purity ($\mathcal F_3=N^2\pi_3$), and put
\[
 F_\lambda(z)=N(z)\bigl(\mathcal F_3(z)-\lambda N(z)^2\bigr),\qquad
 \lambda_{\mathrm{SOS},6}
 =\max\{\lambda:F_\lambda\text{ is a sum of squares of real cubics}\}.
\]
The polynomials are in all $256$ real amplitude coordinates. No
restriction to the conference family or to covariant subsets is imposed.

\begin{theorem}[Complete real cubic-SOS obstruction; TCB]
\label{thm:sextic-sos-obstruction}
The maximal homogeneous degree-six real SOS lower bound satisfies
\[
 \boxed{\quad
 \lambda_{\mathrm{SOS},6}=L_3=\frac{125}{978}<\frac{121}{917}.\quad}
\]
In particular, for every $\lambda>125/978$ there
is no identity
\[
 F_\lambda(z)=\sum_a|p_a(z)|^2+\sum_b|q_b(z,\bar z)|^2
\]
with fixed-coefficient holomorphic cubics $p_a$ and bidegree-$(2,1)$
cubics $q_b$, nor a real homogeneous cubic SOS.
This includes arbitrary local $[2,1]$ sectors and cross-focused
couplings. It excludes neither higher-degree SOS nor non-SOS
inequalities or state-dependent reductions.
\end{theorem}
\begin{proof}
Every squared holomorphic cubic is a positive linear functional of
$|z\rangle\langle z|^{\otimes3}$. Every squared mixed cubic is a
positive linear functional of
$|zz\widetilde z\rangle\langle zz\widetilde z|$, since spin flip
is invertible on the conjugate factor. Average any proposed real
cubic SOS over global phase: only opposite phase weights pair, so
its terms become squared moduli of their $(3,0)$ and $(2,1)$
components, since $(\operatorname{Re}p)^2+(\operatorname{Im}p)^2=|p|^2$. Twirling their positive Gram matrices over the local
unitaries, sites and conjugation preserves positivity and $F_\lambda$; an average of sums of squares is
again a sum of squares.
By Theorem~\ref{thm:spinflip-ppt-hierarchy}, the complete three-copy
Fourier blocks contain both kinds of Gram forms, with no omitted
local representation sector.

Polynomial equality on every $z$ implies equality of the compressed
Hermitian coefficient operators on $\operatorname{Sym}^3H$: the
expectations in all $z^{\otimes3}$ determine such an operator by
its polynomial coefficients. Thus the twirled identity holds on
any feasible averaged moment, not only on physical pure powers.
The argument uses three properties of the moment: its twirled operator $X$ is positive semidefinite on every
$\omega$-block of the $\operatorname{Sym}^3H$ sector, the flipped image $\mathcal T(X)$ is positive semidefinite on
every $\gamma$-block of the $\operatorname{Sym}^2H\otimes H$ sector, and its word values satisfy the one-site
antisymmetriser relations, so that the pairing is well defined. In the isotypic decomposition of the local
three-copy algebra the compressed blocks $\omega_1,\dots,\omega_7$ and $\gamma_0,\dots,\gamma_7$ have sizes $(1,1,1,3,5,11,21)$ and
$(1,1,2,4,8,16,32,64)$.
The boundary pseudo-moment $m^*/70416$ of Theorem~\ref{thm:threecopy} has positive semidefinite blocks,
satisfies the relations and has objective value $125/978$, so it evaluates the left side as
\[
 F_\lambda(m^*/70416)=\frac{125}{978}-\lambda,
\]
while every positive Gram form pairs nonnegatively with it. An SOS representation therefore forces
$\lambda\le125/978$.
Conversely, the TCB dual certificate gives an SOS at $125/978$: its positive
$\omega$ and $\gamma$ Gram forms factor into cubic squared moduli,
and its antisymmetriser relations vanish identically.
\end{proof}

\noindent\emph{Remark.} The relations have exact rank $23$ (TCB); the $37-23=14$ free coordinates match
the independent invariant count from~\eqref{eq:fqe-invariant-dimension} (HI). This count is not
used in the proof.
